\documentclass[a4paper,landscape,fontset=fandol]{ctexart}
\usepackage[landscape,margin=6.5mm,top=7mm,bottom=7mm]{geometry}
\usepackage{amsmath,amssymb,mathtools,bm}
\usepackage{multicol}
\usepackage[most]{tcolorbox}
\usepackage{xcolor}
\usepackage{enumitem}
\usepackage{booktabs,tabularx,array}
\usepackage{fancyhdr}
\usepackage{microtype}
\usepackage{titlesec}

\definecolor{navy}{HTML}{17324D}
\definecolor{teal}{HTML}{087E8B}
\definecolor{sky}{HTML}{EAF5F7}
\definecolor{gold}{HTML}{C47A00}
\definecolor{cream}{HTML}{FFF6E5}
\definecolor{red}{HTML}{B33A3A}
\definecolor{ink}{HTML}{17212B}
\definecolor{muted}{HTML}{52616B}

\color{ink}
\setlength{\parindent}{0pt}
\setlength{\parskip}{1.1pt}
\setlength{\columnsep}{3.2mm}
\setlength{\columnseprule}{0.25pt}
\renewcommand{\columnseprulecolor}{\color{teal!35}}
\setlist[itemize]{leftmargin=3.2mm,itemsep=0.2pt,topsep=0.5pt,parsep=0pt}
\setlist[enumerate]{leftmargin=3.5mm,itemsep=0.2pt,topsep=0.5pt,parsep=0pt}
\pagestyle{fancy}
\fancyhf{}
\renewcommand{\headrulewidth}{0pt}
\renewcommand{\footrulewidth}{0pt}
\fancyfoot[L]{\scriptsize COMP3670 / COMP6670 · Lectures 1--6b}
\fancyfoot[C]{\scriptsize 双面打印：A4 横向，短边翻转}
\fancyfoot[R]{\scriptsize \thepage/2}

\newtcolorbox{boxx}[2][]{
  enhanced,breakable=false,colback=white,colframe=teal!60!navy,
  boxrule=0.45pt,arc=1.2mm,left=1.2mm,right=1.2mm,top=0.8mm,bottom=0.8mm,
  before skip=1.1mm,after skip=1.1mm,
  title={#2},fonttitle=\bfseries\color{white},colbacktitle=navy,
  boxed title style={arc=0.8mm},attach boxed title to top left={xshift=0.7mm,yshift=-0.5mm},
  #1
}
\newtcolorbox{keybox}[1]{
  enhanced,colback=cream,colframe=gold!80!black,boxrule=0.5pt,arc=1mm,
  left=1.2mm,right=1.2mm,top=0.8mm,bottom=0.8mm,before skip=1mm,after skip=1mm,
  title={#1},fonttitle=\bfseries\color{white},colbacktitle=gold!90!black
}
\newcommand{\term}[1]{\textbf{\color{teal!80!black}#1}}
\newcommand{\warn}[1]{\textbf{\color{red}#1}}
\newcommand{\sense}[1]{\textcolor{muted}{\ensuremath{\Rightarrow} #1}}
\newcommand{\E}{\mathbb{E}}
\newcommand{\R}{\mathbb{R}}
\newcommand{\ind}{\mathbf{1}}
\newcommand{\argmin}{\operatorname*{arg\,min}}
\newcommand{\argmax}{\operatorname*{arg\,max}}
\newcommand{\rank}{\operatorname{rank}}
\newcommand{\nullity}{\operatorname{nullity}}
\newcommand{\Span}{\operatorname{span}}

\begin{document}
\fontsize{6.70}{7.50}\selectfont

% ========================= FRONT =========================
\begin{center}
  {\fontsize{14}{15}\selectfont\bfseries\color{navy} COMP3670 概念速查表}\quad
  {\large\bfseries\color{teal}正面：学习、概率、泛化与模型}
  \hfill {\scriptsize 公式不是终点：先问“对象是什么、平均在哪里、复杂度由谁控制？”}
\end{center}
\vspace{-1.5mm}

\begin{multicols*}{5}

\begin{boxx}{1. 机器学习到底在做什么？}
\term{归纳学习}：从有限样本发现规律，并希望它在未见数据上仍成立。算法只是“从候选规则中选一个”的工具；真正的模型还包括数据假设、假设集合和损失。

\term{三类学习}
\begin{itemize}
  \item 监督学习：观测 $(x_i,y_i)$；分类的 $Y$ 离散，回归的 $Y$ 连续。
  \item 无监督学习：只有 $x_i$；寻找聚类、低维结构或密度。
  \item 强化学习：状态 $\to$ 动作 $\to$ 奖励；学习最大化累计奖励的策略。
\end{itemize}

监督学习的四个核心对象：
\[
\begin{gathered}
D=\{(x_i,y_i)\}_{i=1}^N,\quad f^*:\mathcal X\to\mathcal Y,\\
\mathcal H=\{h\},\quad \mathcal A:D\mapsto \hat h
\end{gathered}
\]
\term{目标函数 $f^*$} 是现象允许的最佳预测；\term{假设集 $\mathcal H$} 是学习前允许考虑的规则；\term{学习算法 $\mathcal A$} 用数据选出 $\hat h$。

\warn{关键区分}：训练误差小只是“解释样本”，泛化误差小才是“学到现象”。能拟合一切的模型不可证伪，也最容易过拟合。
\end{boxx}

\begin{boxx}{2. 损失、经验风险与真正风险}
\term{损失} $\ell((x,y),h)$ 衡量一次预测的代价。
\begin{itemize}
  \item 二分类：$\ell=\ind\{h(x)\neq y\}$。
  \item 回归：$\ell=(h(x)-y)^2$。
  \item 概率预测：交叉熵 $-y\log\eta(x)-(1-y)\log(1-\eta(x))$。
\end{itemize}
\[
\begin{aligned}
E_{\rm in}(h)&=\frac1N\sum_{i=1}^N\ell((x_i,y_i),h),\\
E_{\rm out}(h)&=\E_{(X,Y)\sim P}[\ell((X,Y),h)].
\end{aligned}
\]
\sense{$E_{\rm in}$ 是样本平均，可计算；$E_{\rm out}$ 是对未知总体的平均，才是最终目标。}

\term{ERM（经验风险最小化）}
\[
\hat h=\argmin_{h\in\mathcal H}E_{\rm in}(h).
\]
ERM 只负责在 $\mathcal H$ 内找训练表现最好的规则；它不自动保证泛化。
\end{boxx}

\columnbreak

\begin{boxx}{3. 概率语言：随机性来自哪里？}
\term{随机变量}把随机结果映成数；分布描述可能取值及其概率。离散型用概率质量，连续型用密度（点概率通常为 0，区间概率是积分）。
\[
\begin{aligned}
\E[X]&=\sum_x xP(X=x)\quad(\text{离散}),\\
\E[X]&=\int x p(x)\,dx\quad(\text{连续}),\\
\operatorname{Var}(X)&=\E[(X-\E X)^2].
\end{aligned}
\]
\term{条件分布} $P(Y\mid X=x)$：知道特征后，对标签还剩多少不确定性。\term{独立}表示联合分布可分解；i.i.d. 样本意味着每个观测独立且来自同一 $P$。

\warn{不要混淆}：$f^*$ 并不必然把每个样本预测正确；噪声存在时，同一个 $x$ 可对应不同 $y$。
\end{boxx}

\begin{boxx}{4. Bayes 最优规则：理论下限}
二分类令 $\eta(x)=P(Y=1\mid X=x)$。在两类错误代价相同时，
\[
f^*(x)=\ind\{\eta(x)>1/2\}.
\]
\sense{选择条件概率更大的类别；它逐点最优，因此总体也最优。若错误代价不同，阈值不再是 $1/2$。}

平方损失下，
\[
\begin{aligned}
f^*(x)&=\E[Y\mid X=x],\\
E_{\rm out}(f^*)&=\E[\operatorname{Var}(Y\mid X)].
\end{aligned}
\]
后式是\term{不可约误差}：即使知道真实规律也无法消除的噪声。

\term{Plug-in 思想}：先估计 $\eta$ 或 $\E[Y\mid X=x]$，再代入最优决策规则。它与“直接在 $\mathcal H$ 上最小化误差”的 ERM 是两条不同路线。
\end{boxx}

\columnbreak

\begin{boxx}{5. 为什么能泛化？}
\term{泛化间隙}
\[
\Delta(h)=|E_{\rm in}(h)-E_{\rm out}(h)|.
\]
固定一个 $h$ 时，Hoeffding 型结论给出
\[
P\{\Delta(h)>\epsilon\}\le 2e^{-2N\epsilon^2}.
\]
但学习后的 $\hat h$ 是\emph{看过数据才选的}，必须同时控制整个 $\mathcal H$。有限 $|\mathcal H|=M$ 时用 union bound：
\[
P\{\exists h:\Delta(h)>\epsilon\}\le 2M e^{-2N\epsilon^2}.
\]
\sense{样本越多，间隙缩小；候选规则越多，挑中“偶然好看”的规则越容易。}

无限假设集不能简单数元素，而要数它在 $N$ 个点上能产生多少种标记。
\begin{itemize}
  \item \term{Shatter coefficient $S(\mathcal H,N)$}：在任意 $N$ 点上可产生的最大二分方式数。
  \item \term{VC 维 $d_{VC}$}：能被 $\mathcal H$ 打碎（实现全部 $2^N$ 种标记）的最大点数。
  \item $d_{VC}<\infty$ 意味着复杂度可控；典型样本复杂度随 $(d_{VC}+\log(1/\delta))/\epsilon^2$ 增长。
\end{itemize}
\term{PAC}：“大概率（$1-\delta$）近似正确（误差至多 $\epsilon$）”。有限 VC 维与二分类的 PAC 可学习性等价。
\end{boxx}

\begin{boxx}{6. 两类误差与模型选择}
\[
\begin{aligned}
\text{超额误差}&\approx \text{逼近误差}+\text{估计误差}\\
\text{逼近误差}&:\ \mathcal H\text{ 太受限},\\
\text{估计误差}&:\ \mathcal H\text{ 太灵活或数据少}.
\end{aligned}
\]
\term{欠拟合}：高训练误差、高测试误差，通常是模型太简单。\term{过拟合}：训练误差低而测试误差高，通常是复杂度相对样本量过大。

\term{正确的数据分工}
\begin{itemize}
  \item 训练集：拟合参数；验证集/交叉验证：选择复杂度、超参数和模型。
  \item 测试集：最后一次、无偏地报告泛化表现；看完再改模型就已变成验证集。
  \item $K$ 折交叉验证：轮流用一折验证、其余训练，平均验证误差；最终用选定设置在全部非测试数据上重训。
\end{itemize}
\end{boxx}

\columnbreak

\begin{boxx}{7. 参数模型：少量参数概括数据}
\term{参数模型} $\{p(\cdot\mid\theta):\theta\in\Theta\subseteq\R^m\}$ 假定分布形状，只估计有限个参数。\term{似然}把观测数据固定、把参数当变量：
\[
\begin{aligned}
L(\theta)&=\prod_{i=1}^N p(z_i\mid\theta),\\
\hat\theta_{ML}&=\argmax_\theta \log L(\theta).
\end{aligned}
\]

\term{逻辑回归}令 $\eta(x)=\sigma(w^Tx+w_0)$，$\sigma(s)=1/(1+e^{-s})$。最大似然等价于最小化交叉熵；以 $1/2$ 为阈值后，边界仍是 $w^Tx+w_0=0$。分数是 log-odds：
\[
\log\frac{\eta(x)}{1-\eta(x)}=w^Tx+w_0.
\]
系数 $w_j$ 每增加 1 单位特征，使 odds 乘 $e^{w_j}$。

\term{感知机} $h_w(x)=\operatorname{sign}(w^Tx+w_0)$；错分 $(x_i,y_i)$ 时更新 $w\leftarrow w+y_ix_i$、$w_0\leftarrow w_0+y_i$。仅在线性可分时保证有限步收敛。
\end{boxx}

\begin{boxx}{8. 非参数分类：用“附近”投票}
统一形式：
\[
\begin{gathered}
\hat\eta(x)=\sum_i W_i(x)\ind\{y_i=1\},\\
W_i\ge0,\qquad \sum_iW_i=1.
\end{gathered}
\]
\begin{tabularx}{\linewidth}{@{}>{\bfseries}lXXX@{}}
\toprule
规则 & 谁投票 & 平滑旋钮 & 小旋钮时\\\midrule
直方图 & 同一格 & 格宽 $l$ & 过拟合\\
$k$NN & 最近 $k$ 个 & $k$ & 过拟合\\
核方法 & 所有人，随距离衰减 & 带宽 $\sigma$ & 过拟合\\\bottomrule
\end{tabularx}
它们依赖\term{空间一致性}：近点标签更相似。特征尺度错误、无关特征多或维度很高时，“近”会失去意义（维数灾难）。训练几乎免费，预测较贵；平滑量必须用验证选择。
\end{boxx}

\columnbreak

\begin{boxx}{9. 树、森林与回归}
\term{决策树/CART}：反复用“$x_j\le\alpha$？”切分空间；选择使加权子节点\term{不纯度}下降最大的切分。常用
\[
\begin{aligned}
I_{\rm Gini}&=2p(1-p),\\
I_{\rm ent}&=-p\log p-(1-p)\log(1-p).
\end{aligned}
\]
树易解释、无需缩放，但边界呈阶梯且对样本非常不稳定；深度、叶数、最小叶样本数控制复杂度。贪心长树\warn{不是}全局 ERM。

\term{Bagging}：bootstrap 重采样，拟合多棵树并投票，主要降低方差。\term{随机森林}再让每次切分只看随机特征子集，以降低树之间相关性；牺牲可解释性换稳定性。

\term{回归}：最小二乘 $\min_\theta\sum_i(y_i-h(x_i;\theta))^2$ 就是平方损失 ERM。线性回归是“对参数线性”，可用非线性基函数 $\phi_j(x)$：$h(x)=\sum_j\theta_j\phi_j(x)$。核回归则用邻近 $y_i$ 的加权平均。
\end{boxx}

\begin{keybox}{正面一眼判断}
先问：输出是标签、数值还是概率？再问：是在估计条件分布（plug-in），还是直接优化经验准则（ERM/贪心）？最后找复杂度旋钮：$|\mathcal H|$、VC 维、$k/\sigma/l$、树深、基函数阶数。
\end{keybox}

\end{multicols*}
\newpage
\fontsize{6.70}{7.50}\selectfont

% ========================= BACK =========================
\begin{center}
  {\fontsize{14}{15}\selectfont\bfseries\color{navy} COMP3670 概念速查表}\quad
  {\large\bfseries\color{teal}反面：向量、聚类与线性代数}
  \hfill {\scriptsize 形状先行：每一步先写维度，再做运算。}
\end{center}
\vspace{-1.5mm}

\begin{multicols*}{5}

\begin{boxx}{10. 向量：数据点的表示}
$x=(x_1,\ldots,x_d)^T\in\R^d$ 是有顺序的特征列表；顺序和单位都属于表示的一部分。
\[
\begin{aligned}
x+y&=(x_j+y_j)_j,\qquad ax=(ax_j)_j,\\
x^Ty&=\sum_jx_jy_j.
\end{aligned}
\]
\term{范数}衡量“大小”，并诱导距离 $d(x,y)=\|x-y\|$：
\[
\begin{aligned}
\|x\|_1&=\sum_j|x_j|,\qquad \|x\|_\infty=\max_j|x_j|,\\
\|x\|_2&=\sqrt{\sum_jx_j^2}.
\end{aligned}
\]
欧氏内积给出 $x^Ty=\|x\|_2\|y\|_2\cos\theta$；正交即 $x^Ty=0$。

\warn{尺度警告}：距离会被数值范围大的特征支配。距离模型前通常要标准化；类别编码也不能随意当作有序距离。
\end{boxx}

\begin{boxx}{11. K-Means：用质心压缩数据}
给定 $K$ 个簇，交替：
\begin{enumerate}
  \item \term{分配}：每点归入最近质心；
  \item \term{更新}：每个质心改为簇内均值。
\end{enumerate}
目标是最小化簇内平方和（SSE）：
\[
\begin{aligned}
J&=\sum_{k=1}^K\sum_{x_i\in C_k}\|x_i-\mu_k\|_2^2,\\
\mu_k&=\frac1{|C_k|}\sum_{x_i\in C_k}x_i.
\end{aligned}
\]
每一步都不增大 $J$，有限种分配意味着算法终止；但只保证\term{局部最优}，结果依赖初始化。常用多次随机重启或 k-means++。

适合近似球状、方差和大小相近的簇；对离群点、非凸簇、尺度差异敏感。$K$ 是复杂度选择，不应只看训练 SSE（$K$ 增大时它必降）。
\end{boxx}

\columnbreak

\begin{boxx}{12. 层次聚类：合并历史也是结果}
凝聚式算法从单点簇开始，每步合并距离最近的两簇，直到只剩一个。\term{linkage} 定义簇间距离：
\[
\begin{aligned}
d_{\rm single}(A,B)&=\min_{a\in A,b\in B}d(a,b),\\
d_{\rm complete}(A,B)&=\max_{a\in A,b\in B}d(a,b).
\end{aligned}
\]
single linkage 能找细长/非凸结构，但会“链化”；complete linkage 偏好紧凑、直径小的簇，对离群点敏感。

\term{树状图 dendrogram} 的高度是合并距离；在某高度横切得到聚类。它保留多尺度结构，无需预先固定 $K$，但距离与 linkage 的选择会改变答案。
\end{boxx}

\begin{boxx}{13. 矩阵与线性映射}
$A\in\R^{m\times d}$ 有 $m$ 行、$d$ 列。矩阵代表线性映射 $T:\R^d\to\R^m$：
\[
\begin{gathered}
T(x)=Ax,\\ A(\alpha x+\beta y)=\alpha Ax+\beta Ay.
\end{gathered}
\]
两个互补视角：
\[
\begin{aligned}
Ax&=\sum_{j=1}^d x_j a_j &&(\text{列的线性组合}),\\
(Ax)_i&=r_i^Tx &&(\text{每行做一次内积}).
\end{aligned}
\]
\sense{列告诉你输出能由哪些方向组合；行告诉你每个输出坐标如何读取输入。}

形状规则：$A_{m\times d}B_{d\times p}=C_{m\times p}$，内维必须相等。$AB$ 表示\term{先做 $B$，后做 $A$}；通常 $AB\neq BA$。
\[
\begin{gathered}
A(BC)=(AB)C,\\ A(B+C)=AB+AC,\\ AI=IA=A.
\end{gathered}
\]
\end{boxx}

\columnbreak

\begin{boxx}{14. 转置、批处理与二次型}
$A^T$ 交换行列：$(A^T)_{ij}=A_{ji}$。
\[
\begin{gathered}
(A^T)^T=A,\qquad (A+B)^T=A^T+B^T,\\
(AB)^T=B^TA^T.
\end{gathered}
\]
乘积转置\warn{顺序反转}。若数据矩阵 $X\in\R^{N\times d}$ 每行是一个样本，则同时线性变换所有样本可写为 $XA^T$（取决于列/行样本约定，先验形状）。

最小二乘的矩阵形式：
\[
\begin{aligned}
L(w)&=\|Xw-y\|_2^2\\
&=(Xw-y)^T(Xw-y).
\end{aligned}
\]
展开时标量满足 $y^TXw=(Xw)^Ty$，故交叉项合成 $-2y^TXw$。

\term{两层纯线性网络不会增加表达力}：$A(Bx)=(AB)x$。若没有非线性激活，多层仍只是一个线性映射。
\end{boxx}

\begin{boxx}{15. 线性方程组与消元}
$Ax=b$ 等价于同时满足所有行方程，也等价于“$b$ 是否能由 $A$ 的列线性组合得到”。增广矩阵写作 $[A\mid b]$。

三种初等行操作（换行、非零倍缩放、一行加另一行倍数）都可逆，因此不改变解集。Gaussian elimination 得 REF；继续把主元化为 1 且清空主元列得到唯一的 RREF。

\term{主元变量}由方程决定；\term{自由变量}可任取参数。看 RREF 快速判断：
\begin{itemize}
  \item 出现 $[0\ \cdots\ 0\mid c]$ 且 $c\neq0$：无解；
  \item 一致且每个变量列都有主元：唯一解；
  \item 一致且有自由变量：无穷多解。
\end{itemize}
\warn{行操作用于保留方程解；列操作会改变变量意义，通常不能混用。}
\end{boxx}

\columnbreak

\begin{boxx}{16. 逆矩阵：什么时候变换可撤销？}
方阵 $A$ 可逆若存在 $A^{-1}$ 使 $A^{-1}A=AA^{-1}=I$。求逆可对 $[A\mid I]$ 行化简：若得到 $[I\mid B]$，则 $B=A^{-1}$。

等价判断（$A\in\R^{n\times n}$）：
\begin{itemize}
  \item $A$ 可逆；$Ax=b$ 对每个 $b$ 有唯一解；
  \item RREF$(A)=I$；每行每列都有主元；
  \item 列向量线性无关且张成 $\R^n$；
  \item $\operatorname{Null}(A)=\{0\}$；$\rank(A)=n$。
\end{itemize}
\[
\begin{gathered}
(AB)^{-1}=B^{-1}A^{-1},\\ (A^T)^{-1}=(A^{-1})^T.
\end{gathered}
\]
\term{概念}：逆不是“除以矩阵”，而是撤销线性变换；信息被压扁到更低维时便无法撤销。
\end{boxx}

\begin{boxx}{17. Span、子空间、列空间与零空间}
\term{线性组合} $c_1v_1+\cdots+c_kv_k$；所有组合构成
\[
\Span\{v_1,\ldots,v_k\}.
\]
\term{子空间} $S\subseteq\R^n$ 必须含 $0$，且对加法、数乘封闭。检验常用：任取 $u,v\in S$ 与标量 $a,b$，证明 $au+bv\in S$。

对 $A=[a_1\cdots a_d]\in\R^{m\times d}$：
\[
\begin{aligned}
\operatorname{Col}(A)&=\Span\{a_1,\ldots,a_d\}\subseteq\R^m,\\
\operatorname{Null}(A)&=\{x:Ax=0\}\subseteq\R^d.
\end{aligned}
\]
\sense{$Ax=b$ 有解当且仅当 $b\in\operatorname{Col}(A)$；零空间描述被 $A$ 压成 0 的输入方向。}

\warn{列空间基}要取\emph{原矩阵}的主元列；行化简后的主元列只告诉你索引，因为行操作会改变列空间。
\end{boxx}

\columnbreak

\begin{boxx}{18. 独立、基、维数、秩}
$v_1,\ldots,v_k$ \term{线性无关}若
\[
\begin{gathered}
c_1v_1+\cdots+c_kv_k=0\\ \Rightarrow\ c_1=\cdots=c_k=0.
\end{gathered}
\]
否则某向量可由其他向量组合，信息有冗余。把向量作列组成 $V$，求 $V c=0$：只有零解 $\Leftrightarrow$ 独立。

\term{基} = “张成 + 独立”：既不缺方向，也无冗余。一个子空间的所有基含相同数量向量，称为\term{维数}。

\term{秩}是主元数，也是列空间维数；\term{零度}是自由变量数，也是零空间维数。对 $A\in\R^{m\times d}$：
\[
\begin{aligned}
\rank(A)&=\dim\operatorname{Col}(A),\\
\nullity(A)&=\dim\operatorname{Null}(A),
\end{aligned}
\]
\[
\boxed{\rank(A)+\nullity(A)=d}
\]
上式即\term{秩--零度定理}。
\sense{每个输入方向要么贡献可见输出（秩），要么落入被消除的方向（零度）。}
\end{boxx}

\begin{keybox}{反面解题路线}
向量/矩阵题先写形状；聚类题先确认距离与尺度；方程题做 RREF 看矛盾行和自由变量；空间题把“是否张成”化为 $Ax=b$，把“是否独立”化为 $Ax=0$；最后用主元数和秩--零度定理互相校验。
\end{keybox}

\end{multicols*}
\end{document}
